\section{Related Work}
\label{sec:related}

\paragraph{Anisotropy in contextual embeddings.} \citet{ethayarajh2019contextual} document that BERT, ELMo, and GPT-2 produce contextual embeddings that occupy a narrow cone of the hypersphere at every layer, with the effect intensifying with depth. \citet{mu2018allbutthetop} show that subtracting the mean and removing top principal components (``all-but-the-top'') makes static word embeddings more isotropic and improves downstream tasks. \citet{gao2019representation} identify a parallel phenomenon in language generation models, where weight-tied softmax training pushes learned embeddings into a degenerate cone, and propose a cosine-based regularizer to widen it. For sentence-level encoders, BERT-flow \citep{li2020sentenceembeddings} learns an invertible flow to an isotropic Gaussian, and WhiteningBERT \citep{huang2021whiteningbert} applies a closed-form whitening transform at inference. \citet{rajaee2022isotropy} extend the diagnostic analysis to multilingual BERT and show that anisotropy is present across all languages, with substantial variation between high- and low-resource languages. These works diagnose and mitigate anisotropy at the sentence level using post-hoc or pooling-based fixes. We study token-level geometry in late-interaction retrieval, where pooling does not occur and the regularizer must operate at training time.

\paragraph{Contrastive learning and uniformity.} \citet{wang2020understanding} decompose the asymptotic InfoNCE loss into an alignment term and a uniformity term and show that directly optimizing both recovers contrastive training. Uniformity on the hypersphere is closely related to isotropy: a uniformly distributed set of $L_2$-normalized vectors has zero mean pairwise cosine similarity. SimCSE \citep{gao2021simcse} relies on this connection implicitly, using dropout as data augmentation to induce both alignment and uniformity for sentence embeddings. \citet{xiao2023isotropy} analyze how contrastive losses partially induce isotropy in sentence encoders. Our regularizer is a token-level analog of the Wang and Isola uniformity term, applied inside sequence models where InfoNCE operates at the sentence level but the score function operates at the token level.

\paragraph{Late-interaction retrieval.} ColBERT \citep{khattab2020colbert} introduces late interaction with token-level MaxSim scoring. ColBERTv2 \citep{santhanam2022colbertv2} improves quality and efficiency through denoised supervision and residual compression of token vectors. PLAID \citep{santhanam2022plaid} provides an efficient retrieval engine for compressed ColBERTv2 indexes. Recent multilingual extensions include ColBERT-XM \citep{louis2024colbertxm} and Jina-ColBERT-v2 \citep{jha2024jinacolbert}, both of which scale late interaction to many languages through modular architectures and additional training data. These works improve quality through better training data, supervision, or compression. We improve quality through a geometric prior on the token vectors themselves.

\paragraph{Multilingual dense retrieval.} XLM-R \citep{conneau2020xlmr} is the dominant multilingual backbone for dense retrievers. LaBSE \citep{feng2022labse} combines masked and translation language modeling with margin-based contrastive learning for cross-lingual sentence embeddings. Contriever \citep{izacard2022contriever} introduces unsupervised contrastive pre-training for dense retrieval and extends it to multiple languages. Multilingual E5 \citep{wang2024me5} uses large-scale weakly supervised contrastive pre-training followed by supervised fine-tuning. The standard multilingual retrieval benchmarks include MIRACL \citep{zhang2023miracl} for monolingual ad-hoc retrieval across 18 languages and mMARCO \citep{bonifacio2021mmarco} for multilingual passage ranking. These efforts focus on data and supervision; we hold both fixed and intervene only on training-time geometry.

\paragraph{Isotropy regularization for retrieval and beyond.} The closest prior work to ours applies isotropy techniques to ColBERT or to training-time regularization. \citet{jung2023isotropic} apply post-hoc normalizing flow and whitening to ColBERT and RepBERT, both sentence-level and token-level, on English retrieval benchmarks. \citet{kim2024hil} propose a hybrid isotropy/anisotropy ColBERT architecture for zero-shot transfer on BEIR. Both are monolingual and do not target the cross-lingual regime where pretrained signal asymmetry is the dominant source of geometric collapse. I-STAR \citep{rudman2024istar} adds a differentiable isotropy regularizer to fine-tuning and reports that \emph{decreasing} isotropy can help classification tasks, suggesting that the productive direction is task-dependent. In a separate line, \citet{ji2023isotropic} apply isotropy interventions to zero-shot cross-lingual transfer on classification, not retrieval. Our paper occupies a different point in this design space: training-time, token-level, multilingual, late-interaction, low-resource. We address the I-STAR result directly: for MaxSim retrieval, the productive direction is toward higher isotropy across the full token manifold.
